% ---------------------------------------------------------------------
% CHAPTER 1 — INTRODUCTION
% ---------------------------------------------------------------------
\chapter{Introduction \label{ch:introduction}}

% 1.1 The problem, stated once, concretely.
%   - Architecture drifts/erodes; documentation goes stale; mappings are missing.
%   - SAR and C2A mapping defined and distinguished (recovery = discover the
%     decomposition; mapping = assign entities given partial known architecture).
%   - Cost anchor: expert C2A mapping is ~107 person-hours/system; Chromium took
%     ~2 years. This is the "why automate" evidence, not rhetoric.
\section[Architectural Erosion and Recovery]{Software Architecture Erosion and the Recovery}

% 1.2 Why the obvious approaches are insufficient (motivates the whole thesis).
%   - Dependency-only clustering: "connected by an edge" != "same module".
%   - Hand-tuned fusion (SARIF): fixed weights don't generalize across systems/scale.
%   - The gap: no single signal suffices, and hand-weighting them doesn't either.
\section{Why architectural signal is hard to use}

% 1.3 The thesis statement (the spine). State it plainly:
%   Architectural structure is distributed across several implementation signals
%   (typed dependencies, folder organization, lexical identifiers, member-level
%   detail). Learned graph representations integrate them better than hand-tuned
%   fusion; and recovery is bounded not by the encoder but by whether the target
%   architecture is reflected in the chosen signal and whether the model's reach
%   matches system scale and granularity.
\section{Thesis statement and scope}

% 1.4 Contributions of the thesis as a whole (not per-paper yet).
%   - C1: learned representations beat hand-tuned fusion, most clearly at scale.
%   - C2: isolation of what each signal contributes and where it fails.
%   - C3: the same signals/systems hold across unsupervised recovery AND
%         supervised mapping (the duality that no single paper can claim).
%   - C4: directory structure as a strong architectural proxy, shown by removing it.
\section{Contributions}

% 1.5 Included papers + author roles + a one-line "what each isolates".
\section{Included publications}

% 1.6 Thesis outline (short).
\section{Outline}
% \section{Motivation: Architecture Erosion and the Cost of Technical Debt}
% \label{sec:motivation}

% ---------------------------------------------------------------------
% CHAPTER 2 — BACKGROUND
% ---------------------------------------------------------------------
\chapter{Background \label{ch:background}}

\section[Architectural Drift and Reconstruction]{Architectural Drift, Erosion and Reconstruction \label{sec:bg-architecture}}

Software architecture describes the high-level structure of a system in terms of its principal elements, their responsibilities, and the relations among them, and provides the basis for reasoning about maintainability, security, and other quality attributes~\cite{clements2003}. This structure is often reflected in the implemented system as a modular decomposition, a partition of files into modules that group related functionality~\cite{parnas1972}. Such decompositions matter for program understanding, maintenance, and refactoring, but documented architecture is often missing, incomplete, or quickly outdated as systems evolve~\cite{clements2003}. As a system evolves, developers introduce dependencies that were not anticipated in the original design, refactor code across module boundaries, and add functionality that fits no existing module, and documentation is rarely maintained in step with these changes. The resulting gap between the intended design and the implementation is described using two related terms: architectural drift, where the implementation grows away from the intended design without directly violating it, and architectural erosion, where the implementation directly violates the intended design~\cite{de2012controlling}.
 
Architectural drift and erosion leave the documented architecture unreliable, complicating maintenance and evolution when accurate architectural knowledge is most needed. Two activities respond to this problem. Software Architecture Recovery (SAR) addresses the case where no trusted decomposition is available, by inferring a modular decomposition directly from the source code and its dependencies. Manually organizing hundreds or thousands of files into coherent modules does not scale, and recovering the architecture of even a medium-sized system can require substantial expert effort~\cite{garcia2013}. Code-to-architecture (C2A) mapping applies in the opposite situation: an intended architecture is known, and the task is to assign source files to its modules, which is the enabling step for reflexion modeling~\cite{reflexion1}, where module-level dependencies are classified as convergent, divergent, or absent. Reflexion modeling depends on such a mapping between source code entities and architectural modules, and this mapping is often missing, incomplete, or expensive to maintain manually~\cite{ali2018architecture}.
 
The two problems share a common structure: in both, the unit of assignment is the source file, and the evidence available is drawn from the same implementation artifacts: dependencies, folder structure, and lexical content. Static dependencies capture structural relationships between source-code files. Files that implement related functionality often depend on one another directly, or are connected through short dependency paths, making the dependency graph a useful signal for architecture recovery~\cite{negar}. Dependency types may also carry different architectural meanings; for example, call edges can indicate tighter coupling than type-usage edges~\cite{edge_weights}. Folder structure reflects how developers group files during implementation and is used as evidence by experts recovering ground-truth decompositions~\cite{garcia2013}, but folders are not architectural modules, and the correspondence between the two is partial. Code identifiers encode domain-specific knowledge that often aligns with module membership, but the same vocabulary can recur across modules that share concepts~\cite{sarif}.
 
Prior work in SAR and C2A mapping has approached this distributed signal in two broad ways. One line of work combines several sources by fixing a set of features and weighting them through hand-tuned rules or search over the weight space; the integration is explicit, but the weights are chosen once and do not adapt across systems of different sizes, languages, or design conventions. The other line uses one source at a time, most often the dependency graph, and reasons over it with clustering, search-based optimization, or, more recently, representation learning. Both strategies leave the same question open: how to integrate heterogeneous architectural signals under a single representation that adapts to the system rather than to a fixed configuration. This question is what the remainder of this chapter, and the thesis as a whole, is organized around. Section~\ref{sec:bg-prior} examines how prior approaches have engaged with it, and Section~\ref{sec:bg-gnn} introduces the graph neural network vocabulary that this thesis uses to address it.
 
\section{Prior Approaches}
\label{sec:bg-prior}
 
Recovery and mapping draw on overlapping signals, dependencies, and lexical content most consistently, and both lineages move from fixed, hand-specified integration toward learned integration over time. This section reads them separately, since mapping's shift toward learned models and, more recently, graph representation learning follows the corresponding shift in recovery by several years, but each follows the same underlying trajectory from fixed integration to learned integration.
 
\subsection{Architecture Recovery}
 
Traditional recovery methods treat software artifacts as a dependency graph and recover modules through several related paradigms. A first group formulates recovery as search over a modularization objective: some optimize a single modularity criterion favoring high cohesion within modules and low coupling across them~\cite{bunch_sahc_nahc,dagc,wca}, while others use multi-objective or evolutionary search to balance several criteria at once~\cite{mca_eca}. A second group clusters directly on dependency-derived similarity or information-theoretic criteria~\cite{cct,hcc_2019,limbo}, and a third emphasizes practicality, determinism, or scalability in the clustering procedure itself~\cite{fca,neighborhood_tree_2019,gma_2021}. These methods can be effective when dependencies reflect architectural structure, but recovery quality is sensitive to how the dependency graph is constructed and how dependency types are weighted.
 
Semantic approaches use identifiers and comments instead of, or alongside, dependencies, and can outperform dependency-only methods in some systems~\cite{comparative_survey}. ARC, ZBR, and Risi et al. cluster artifacts in a latent semantic space learned from source code, whether through topic models, TF-IDF zones, or LSI~\cite{arc,zbr,risi2012fold}, and Kargar et al.~\cite{kargar2017semantic} build a semantic dependency graph from lexical cues and optimize it by local search. Anquetil et al.~\cite{file_clustering} instead use file names and directory paths as a proxy for architecture. Non-structural signals of this kind can be noisy in isolation, since reused terminology and utility directories mix entities from several modules.
 
Hybrid approaches combine multiple sources to improve robustness. Weighted similarity fusion with search-based partitioning and many-objective remodularization augment a structural objective with semantic or evolution-related criteria~\cite{sc_misra,escrc,nsga3,sadat2019multi}. Puhala et al.~\cite{ensemble_recovery} combine several clustering results into a consensus solution, and Yang et al.~\cite{evol} use semantic cues to guide label propagation over the dependency graph. Naim et al.~\cite{naim2017cchd} combine structural and contextual information through coordinated clustering, adapting the combination through incremental developer feedback rather than fixed weights. SADE~\cite{sade} and SARIF~\cite{sarif} both fuse signals by building a weighted dependency graph and applying community detection; SARIF, used as a baseline in the recovery studies in this thesis, extends this pipeline with directory similarity and Louvain community detection. Across this group, the fusion is usually a fixed rule or a heuristic weighting rather than a representation learned jointly from the signals.

Representation learning has also been applied to recovery. NEGAR~\cite{negar} learns embeddings from the dependency graph through random walks and clusters the resulting vectors, and reports that recovery quality varies across dependency types, motivating methods able to exploit several relations jointly. NEGAR runs a separate pipeline per dependency type, however, and does not incorporate non-structural signals. These limitations motivate the recovery studies synthesized in this thesis. They examine multi-relational graph representation learning through reconstruction and contrastive learning objectives, the integration of directory and lexical information~\cite{gaer,arl_2026}, and the extent to which the choice of input graph interacts with encoder reach~\cite{structsar}.

More recently, LLMs have reached recovery as well: Rukmono et al.~\cite{rukmono2024deductive} use chain-of-thought prompting with GPT-4 to map implementation units to a predefined reference architecture, which requires an available architectural model and raises scalability concerns as system size grows. SemArc~\cite{sar_sem} instead uses an LLM to interpret code and architectural semantics against canonical patterns, then clusters components with an anchor-based strategy, and related work applies language models to end-to-end recovery from source code~\cite{amalfitano2025automated,archagent}.

\subsection{Mapping}
 
Tzerpos and Holt~\cite{tzerpos1997orphan} introduced orphan adoption as the incremental assignment of unmapped entities to a documented architecture. Christl et al.~\cite{christl2007automated} built on this with a semi-automatic method that recommends mappings through attraction functions and leaves the decision to the architect, and Bittencourt et al.~\cite{bittencourt_IR} studied attraction functions derived from information retrieval, assigning each entity to the module its text most resembles. These methods rely on the same two signals as recovery, lexical similarity and dependency proximity to already-mapped entities, but combine them through fixed, manually specified rules, the mapping analog of the fusion schemes above.
 
A later shift replaces these hand-designed rules with learned models. Olsson et al.~\cite{olsson2022automatically} propose a semi-supervised Naive Bayes mapper that combines lexical and dependency-derived features through Concrete Dependency Abstraction, reporting that 15--20\% of entities must be mapped beforehand to reach acceptable performance. Florean et al.~\cite{florean2021textclassifierscomparison} compare text classifiers for the same task and show that both the classifier and the text representation affect mapping quality. InMap~\cite{InMap} instead ranks candidate mappings interactively from natural language module descriptions as a recommendation system, and requires no pre-mapped entities at all. In these methods, dependency information, where used, typically enters as a feature attached to individual entities rather than as an explicit graph structure. Jabeen et al.~\cite{jabeen2025graph} took the corresponding step to graph representation learning, evaluating GCN-based models on homogeneous and dependency-type-aware graphs and showing that a small labeled set can be expanded through self-training.
 
LLMs have also been applied to mapping: Johansson et al.~\cite{johansson2024mapping} evaluate LLM-based embeddings alongside TF-IDF and GloVe as input to method-level classifiers, and in follow-up work find that classification performance depends less on lexical overlap or class balance and more on how separable the classes are in the embedding space, an effect that varies with documentation granularity~\cite{johansson2026mapping}. These models exploit lexical and semantic signal at a scale earlier methods could not, but typed dependencies stay implicit in their input.

Across both recovery and mapping, a common methodological limitation remains. Existing approaches either rely on a single source of information, combine several through predefined weights, or, where representation learning is applied, learn from dependencies alone. To our knowledge, no prior approach jointly learns a representation from dependencies, directory structure, and lexical content, nor systematically isolates the contribution of each source. The studies comprising this thesis address these limitations. The following section introduces graph neural networks as the technical framework used to represent and learn from these heterogeneous sources of software-system information.
 
\section{Graph Neural Networks}
\label{sec:bg-gnn}
 
Graph neural networks (GNNs) provide a way to learn from several architectural signals within a common representation learning framework. A GNN encoder maps each node, together with its features and graph neighborhood, to a vector representation through message passing~\cite{neuralMP}. At each layer, every node collects the current representations of its neighbors, combines them with its own, and passes the result through a learned transformation. Stacking $L$ layers extends the effective neighborhood from which each node's representation is computed to nodes within $L$-hops. The specific encoder is defined by how neighbors are combined at each layer, and this thesis uses three variants.
 
The graph convolutional network (GCN)~\cite{gcn} aggregates neighboring representations through degree-normalized message passing, which prevents high-degree nodes from dominating solely because they have many dependencies. The graph attention network (GAT)~\cite{gat} instead learns attention weights over neighbors, allowing their contributions to differ according to their representations. Both encoders can be applied either to a graph in which dependency types are collapsed or to a multi-relational graph with relation-specific parameters. When the graph contains multiple node types and typed relations, this thesis also uses the heterogeneous graph transformer (HGT)~\cite{hgt}, which parameterizes attention by meta-relation. Thus, the contribution of a neighbor depends jointly on the source-node type, target-node type, and edge type.
 
Two properties of graphs determine how well message-passing encoders can support downstream tasks, and both recur throughout the thesis: homophily and reach. A graph is homophilic with respect to a labeling if nodes that share a label tend to be connected. Message passing accumulates evidence from neighbors, so under homophily the aggregated representation is a coherent summary of the node's local context, and under heterophily the same aggregation mixes evidence from nodes with different labels and blurs the distinction that the downstream task needs to preserve~\cite{gnn_homophily}. For architecture recovery over a dependency graph, homophily corresponds to dependencies being concentrated within architectural modules rather than between them. It therefore has a close connection to the architectural objective of high within-module cohesion and low between-module coupling: a decomposition with this structure induces a more homophilic dependency graph with respect to its module assignments. This condition varies across systems and constrains how effectively dependency-driven encoders can recover architectural structure, regardless of their expressive capacity.
 
The second property is the reach of the encoder relative to the graph: the extent of the neighborhood from which a node representation can effectively incorporate information. Nominally, an $L$-layer GNN has an $L$-hop receptive field. Increasing the number of layers, however, does not simply extend useful reach: repeated aggregation mixes distant and local information, causing node representations to become progressively less distinguishable, an effect known as over-smoothing~\cite{gnn_oversmoothing}. In practice, GCN and GAT encoders are typically used with only a small number of layers, limiting the contextual range available to each node. Where the architectural signal relevant to a task is not concentrated within a few hops, this limitation constrains what such encoders can recover. Reach is therefore a property of the encoder together with the graph, rather than of encoder capacity or attention mechanism alone, and it recurs throughout the studies in this thesis as a factor that can compete with, and sometimes outweigh, the choice of encoder architecture.
 
Two further design elements recur throughout the thesis but are introduced in detail where they are applied. Graph autoencoders and their variational form (GAE and VGAE)~\cite{gae} train an encoder to reconstruct observed graph structure from learned node representations, and are used in the unsupervised recovery studies. Contrastive objectives, in particular deep graph infomax (DGI)~\cite{dgi}, train an encoder by distinguishing representations derived from the observed graph from those derived from a corrupted view, and are used where task labels are unavailable but structural regularities can be exploited. Both are training objectives rather than encoder architectures.
\section[Evaluation and Ground Truth]{Evaluation and Ground Truth}
\label{sec:bg-eval}
 
The studies in this thesis evaluate their outputs against expert-produced reference decompositions. Since SAR and C2A mapping produce different outputs, they require different evaluation metrics: recovery is evaluated through partition-agreement metrics, whereas mapping is evaluated through conventional classification metrics. This section introduces these measures and discusses how the reference decomposition should be understood as ground truth. The provenance of the specific reference mappings used in this work is discussed in Chapter~\ref{ch:methodology}.
 
Multiple reasonable architectural decompositions may exist for the same software system, since experts can disagree on architectural boundaries and design intent~\cite{garcia2013}, so a recovered decomposition may differ from the particular reference used for evaluation without necessarily being wrong. No single metric can settle this on its own, and the studies in this thesis instead report a set of complementary metrics for recovery: MoJoFM, A2A, and \cc{} as measures of alignment with the reference, and TurboMQ as a reference-independent measure of modularization quality. The adjusted Rand index~\cite{ari} is reported alongside them as a clustering-agreement measure that corrects for chance and remains meaningful even when the recovered and reference partitions have different numbers of modules.

MoJoFM~\cite{mojofm} is widely used in architecture recovery. It counts the minimum number of move and join operations needed to transform the recovered partition into the reference partition, normalized by the maximum possible number of such operations. Because its join operations are comparatively cheap, MoJoFM is less sensitive to over-fragmentation: many small clusters can still score well when they broadly align with the reference. A2A~\cite{a2a} extends MoJoFM with costs for adding, removing, and moving entities and clusters. It can assign a high score to two architectures that are structurally quite different, overestimating their similarity~\cite{a2a_limitation}. \cc{}~\cite{c2c} measures the percentage of reference modules that are sufficiently covered by a recovered module, where coverage is judged against a threshold on the overlap between the two; conventional thresholds of 0.10, 0.33, and 0.50 correspond to weak, moderate, and majority matches. This makes it sensitive to fragmentation: a reference module split across several recovered modules, with no single one holding a majority share, fails the threshold even where individual file assignments look reasonable under a different metric, and its dependence on a fixed threshold makes it more volatile than the other metrics. The adjusted Rand index (ARI)~\cite{ari} is a widely used clustering-agreement measure, ranging from $-1$ to $1$, that corrects for the level of agreement expected by chance and remains meaningful even when the partitions have different numbers of modules, and, together with MoJoFM, it is the more stable and discriminative of the four. A2A's tendency to overestimate similarity makes it the least discriminative in practice.

Unlike the metrics above, which measure alignment with the reference decomposition, TurboMQ~\cite{turbomq} measures modularity: it rewards intra-cluster cohesion and penalizes inter-cluster coupling, computed directly on the dependency graph without reference to any external partition. TurboMQ can favor decompositions with many clusters, so this thesis reports the normalized TurboMQ (\tmq{}) obtained by dividing by the number of clusters in the recovered decomposition to reduce this bias. A high TurboMQ with a low MoJoFM or ARI means the recovery is coherent in dependency terms but does not match ground truth.
 
For mapping, the evaluation is a per-file classification into known modules, and the standard metrics apply directly. Precision, recall, and their harmonic mean F1 are reported per module and aggregated at the macro and micro level. The macro-average treats every module equally and therefore surfaces performance on small modules that would be hidden by a micro-average dominated by large ones. The mapping studies report both, as the difference between them is itself informative about how the model handles the class imbalance that reference decompositions typically exhibit.
 
This thesis draws its reference decompositions from published, expert-validated benchmarks, so that results remain comparable with prior work and no new judgment is introduced into the evaluation. Matching the reference on every system does not make a method right in any stronger sense than the reference itself.
% ---------------------------------------------------------------------
% CHAPTER 3 — RESEARCH PROBLEM AND QUESTIONS
% ---------------------------------------------------------------------
\chapter{Research Problem and Questions}

\section{Problem Formulation}
\label{sec:problem-formulation}

This thesis considers two related software architecture analysis problems: software architecture recovery (SAR) and source code-to-architecture (C2A) mapping. Both concern assigning source code files to architectural modules. We represent a software system as a graph $\mathcal{G} = (\mathcal{V}, \mathcal{E})$, where nodes represent software artifacts and edges represent relations between them. Depending on the available system information, $\mathcal{G}$ may be homogeneous, with a single node and edge type, or heterogeneous, with multiple node and edge types that capture complementary implementation and organizational relations. Given $\mathcal{G}$, the common objective is to learn representations that support the assignment of each source code file to an architectural module. Both problems are studied at file-level granularity, as source code files are the primary unit of implementation considered in this thesis.

SAR and C2A mapping differ in the information available about the target decomposition. In SAR, no reliable decomposition is assumed to be available. The task is to infer a partition of the files into modules directly from $\mathcal{G}$. SAR is therefore formulated as an unsupervised representation-learning problem and applies when architectural documentation is absent or cannot be trusted. In C2A mapping, the target module set is known in advance, and a subset of files is labeled with its intended module. The task is to assign the remaining files to one of the known modules, making C2A mapping a partially supervised representation-learning problem. It applies when an intended architecture exists, and new, modified, or unclassified files must be aligned with it.

Despite this difference in supervision, both problems are shaped by the same design choices: the information represented in $\mathcal{G}$, the features used to initialize node representations, and the encoder used to learn from the graph. This shared design space enables SAR and C2A mapping to be studied within a common graph-based framework.


\section{Research objectives}
\label{sec:research-objectives}

This thesis is guided by a single overarching research question.

\noindent\textit{How can graph-based representation learning be effectively designed and adapted to support software architecture recovery and source code-to-architecture mapping, and what governs the effectiveness of different graph representations and input signals across these two tasks?}

Three sub-questions break down this overarching question. RQ1 examines the design of graph-based representation learning, including graph representation, relation modeling, encoder architecture, and learning objective. RQ2 examines the architectural information carried by different input signals and the system properties that condition their usefulness. RQ3 examines how the contribution of these signals changes when they are combined within richer graph representations. The comparison across recovery and mapping is not fully symmetric. Both settings examine file-level dependency graphs, multi-relational representations, random-walk and neural message-passing approaches, and GCN/GAT-based encoders. The mapping studies extend this design space to a heterogeneous graph containing folders, files, and member-level artifacts, processed by an HGT-based encoder. The sub-questions are therefore addressed across complementary studies rather than through one common experimental design applied identically to both tasks.

\begin{GrayBox}{RQ1: Graph-based representation learning design.}
How should graph-based representation learning methods be designed to support software architecture recovery and source code-to-architecture mapping, and which choices of graph granularity, relation modeling, encoder, and learning objective are justified by measurable benefit?
\end{GrayBox}

RQ1 examines how graph construction and representation-learning choices affect architectural analysis. In recovery, the methods learn file representations through self-supervised reconstruction or contrastive objectives and cluster them into modules. In mapping, they use semi-supervised classification from an initial set of file-to-module assignments. Across both settings, the question assesses whether different graph representations, encoders, and learning objectives provide a measurable benefit over simpler alternatives.

\begin{GrayBox}{RQ2: Architectural information in input signals.}
What architectural information is carried by dependency structure, project organization, code identifiers, and under which system properties is each source reliable for recovery and mapping tasks?
\end{GrayBox}

This question combines studies that examine dependency-only and organization-augmented graphs, lexical and structural representations, and member-level interaction information. It considers the conditions under which each source supports architectural analysis, including dependency cohesion, graph connectivity, structural reach, and the degree to which organizational structure aligns with the reference architecture.


\begin{GrayBox}{RQ3: Interaction between signal representation and graph design.}
How does the usefulness of structural, lexical, and contextual file representations change when they are combined with progressively extended heterogeneous graph structures, and under which system properties do these sources provide complementary rather than redundant architectural evidence?
\end{GrayBox}

This question is addressed through a controlled heterogeneous-graph study that varies file representations and graph structure independently while holding the learning setting fixed. It tests whether conclusions drawn in RQ1 and RQ2 about the usefulness of particular representations and architectural signals continue to hold once these sources are available within the same heterogeneous model. The analysis distinguishes gains that reflect complementary architectural evidence from cases in which the value of one signal saturates, or where high performance is driven largely by a source such as folder organization already aligning closely with the reference decomposition.

The three RQs are addressed across the two parts of this thesis. Part~I studies unsupervised recovery and contributes to RQ1 through encoder and objective comparisons on multi-relational and heterogeneous software graphs, and to RQ2 through a controlled comparison of dependency-only and folder-augmented representations. Part~II studies semi-supervised mapping and extends the RQ1 evidence to relation-specific encoders and a graph-free neural baseline under identical training, and to a heterogeneous graph transformer over folders, files, and members. RQ3 is addressed principally in the capstone study of Part~II, which isolates the contribution of each signal within a single heterogeneous model and traces how these contributions change as graph structure is progressively extended.
% ---------------------------------------------------------------------
% CHAPTER 4 — METHODOLOGY
% ---------------------------------------------------------------------
\chapter{Method}
\label{ch:methodology}

This chapter presents the research methodology used to address the two problems this thesis studies, software architecture recovery and code-to-architecture mapping, through the graph-based representation-learning artifacts developed across its studies and the reasoning that connects them into one coherent design. It states the paradigm the studies share, how their progression from recovery to mapping was chosen rather than incidental, the unit and graph representations common to all of them, the taxonomy of input signals that recurs across studies, the two learning settings and what distinguishes the questions they can answer, the controlled-comparison logic used to decide when a design choice is justified, the systems and evaluation protocol the studies share, and the threats that this design addresses.

\section{Research Design}
\label{sec:research_design}
The studies in this thesis form an incremental investigation of graph-based representation learning for architectural analysis. Together, they examine how graph structure, relation modeling, input representations, and learning settings affect software architecture recovery and code-to-architecture mapping. Each study extends the evidence of the preceding work by investigating a limitation or uncertainty that remained unresolved, thereby broadening the empirical basis for the research questions of the thesis.

The research starts from the premise that architectural information is distributed across static dependencies, project organization, and source-code content. Graph-based representation learning provides a common framework for combining or isolating these sources: dependency relations define the structural context of a file, while node and edge attributes encode additional organizational, lexical, contextual, or interaction-level information. The studies investigate how these sources should be represented and learned from, rather than assuming that one graph construction, feature representation, or encoder is uniformly appropriate.

The prior graph-based representation learning approach to architecture recovery, NEGAR~\cite{negar}, learns file embeddings from dependency graphs through node2vec-style random walks and clusters them into modules. It evaluates dependency types through separate pipelines and does not incorporate folder organization or lexical information. This leaves open whether typed dependencies can be modeled jointly and whether organizational and lexical signals provide complementary evidence for recovery. GAER~\cite{gaer} addresses this question by representing typed dependencies as relation-specific graph structure and combining folder-path and code identifier as node features within a graph autoencoder framework. It further examines how encoder choice (GCN vs.\ GAT) and cluster granularity affect the result. Because these signals are fused within the same configuration, however, the contribution of each source cannot be separated from the others.

A subsequent recovery study investigates a different self-supervised learning objective. It uses contrastive learning rather than graph reconstruction and augments file-level dependency edges with attributes describing the member-level interactions that give rise to each dependency. This examines whether interaction type and frequency provide useful architectural information beyond the existence of a file-level dependency~\cite{arl_2026}. A further recovery study then examines dependency structure and project organization more directly. It compares dependency-only file graphs with folder-augmented heterogeneous graphs using both random-walk and graph-attention representation learners. The graph-attention variants are trained through a variational graph autoencoder, and the study additionally evaluates positional encoding and diffusion-based structural augmentation to extend the effective reach of shallow message-passing models~\cite{structsar}.

The remaining two studies carry this design space into automatic mapping of source files to architectural modules, where the module set is known, and a subset of files is labeled. The fourth study formulates mapping as semi-supervised node classification with iterative self-training over a relation-aware graph~\cite{mapping_ext}. It also introduces a graph-free multi-layer perceptron trained with the same features and self-training procedure, providing a comparison absent from the recovery studies between graph-based prediction and feature-only prediction. The difference between the models estimates what the graph-based encoder adds beyond the file-level features, although it does not isolate message passing alone because the models differ in parameterization.

The final study extends this mapping design by introducing a dependency-guided contextual code representation and by expanding the graph to include folder and member nodes. It compares structural, lexical, and contextual representations both without message passing and as inputs to the same heterogeneous graph transformer, across graph variants that add or remove folder and member nodes. Thus, rather than directly validating all earlier results, it tests how these representations and graph structures contribute under one matched semi-supervised mapping setting.